\lecture{4.2}{G-Computation with Discrete Data}
\label{lec:4-2}

\subsection{Empirical Mean Outcome Regression}

When $W$ takes only a small number of possible values, the most direct estimator of
the outcome regression is the sample mean within each $(A,W)$ subgroup. For a fixed
treatment value $a$ and covariate value $w$, define
\[
n_{aw}=\sum_{i=1}^n \ind(A_i=a,W_i=w).
\]
If $n_{aw}>0$, the empirical outcome regression is
\[
\Qbar_n(a,w)
=
\frac{\sum_{i=1}^n Y_i\ind(A_i=a,W_i=w)}
       {\sum_{i=1}^n \ind(A_i=a,W_i=w)}.
\]
This is simply the average observed outcome among people with treatment $A=a$ and
covariates $W=w$.

For example, if $W$ is a binary age indicator, then $\Qbar_n(1,0)$ is the sample mean
outcome among younger treated individuals, while $\Qbar_n(0,1)$ is the sample mean
outcome among older untreated individuals.

\begin{keyidea}
With low-dimensional discrete covariates, empirical g-computation is ordinary
stratification. Estimate the mean outcome in each treatment-covariate stratum, then
average those stratum means using the covariate distribution of the target population.
\end{keyidea}

\subsection{Saturated Regression Models}

A saturated model has enough parameters to reproduce every subgroup mean. With
binary $A$ and binary $W$, there are four treatment-covariate subgroups. A saturated
linear mean model is
\[
\E(Y\given A=a,W=w)=\beta_0+\beta_1 a+\beta_2 w+\beta_3 aw.
\]
The fitted means implied by this model are
\[
\begin{array}{c|cc}
& A=0 & A=1\\
\hline
W=0 & \beta_0 & \beta_0+\beta_1\\
W=1 & \beta_0+\beta_2 & \beta_0+\beta_1+\beta_2+\beta_3 .
\end{array}
\]
Because there are four parameters and four subgroup means, the model can exactly
match the empirical subgroup averages whenever each subgroup is observed. The
interaction term $aw$ is what makes the model saturated. Without it, the model would
force the treatment contrast to be the same at both levels of $W$ on the additive
scale.

\begin{remark}
The word ``saturated'' is about flexibility relative to the chosen covariate
description. If $W$ has many categories or multiple components, a truly saturated
model may require many main effects and interactions.
\end{remark}

\subsection{Estimating $\Prob(W=w)$}

The g-computation parameter also requires the marginal distribution of $W$:
\[
\psi(a)=\sum_w \Qbar(a,w)\Prob(W=w).
\]
For discrete $W$, the empirical estimator is the sample proportion
\[
P_n(W=w)
=
\frac{1}{n}\sum_{i=1}^n \ind(W_i=w).
\]
This is the fraction of the sample in covariate stratum $w$. It estimates the covariate
distribution in the target population represented by the sample.

\subsection{Computing $\E\{\Yone\}$}

The g-computation estimator of $\E\{\Yone\}$ is obtained by setting $a=1$:
\[
\psi_n(1)
=
\sum_w \Qbar_n(1,w)P_n(W=w).
\]
For binary $W$, this becomes
\[
\psi_n(1)
=
\Qbar_n(1,0)P_n(W=0)+\Qbar_n(1,1)P_n(W=1).
\]
The first term is the estimated treated mean outcome among people with $W=0$,
weighted by how common $W=0$ is in the full sample. The second term is the analogous
quantity for $W=1$.

\subsection{Computing $\E\{\Yzero\}$}

Similarly,
\[
\psi_n(0)
=
\sum_w \Qbar_n(0,w)P_n(W=w),
\]
and for binary $W$,
\[
\psi_n(0)
=
\Qbar_n(0,0)P_n(W=0)+\Qbar_n(0,1)P_n(W=1).
\]
This estimates the mean outcome that would have been observed if everyone in the
target population had received treatment level $0$.

\subsection{Computing the ATE}

The estimated average treatment effect is the difference between the two estimated
counterfactual means:
\[
\widehat{\ATE}_{\mathrm{gcomp}}
=
\psi_n(1)-\psi_n(0).
\]
Equivalently,
\[
\widehat{\ATE}_{\mathrm{gcomp}}
=
\sum_w
\left\{
\Qbar_n(1,w)-\Qbar_n(0,w)
\right\}P_n(W=w).
\]
This expression shows that empirical g-computation is a weighted average of
stratum-specific treatment contrasts. The weights are not the covariate distributions
inside the treated and untreated groups separately. They are the covariate distribution
of the full target sample.

\subsection{Equivalent Individual-Level Form}

Because $P_n(W=w)$ is the fraction of individuals whose covariates equal $w$, the
same estimator can be written as
\[
\psi_n(a)
=
\frac{1}{n}\sum_{i=1}^n \Qbar_n(a,W_i).
\]
To see this, start from the subgroup formula:
\[
\sum_w \Qbar_n(a,w)P_n(W=w)
=
\sum_w \Qbar_n(a,w)
\left\{
\frac{1}{n}\sum_{i=1}^n \ind(W_i=w)
\right\}.
\]
Move the finite sum over observations outside:
\[
=
\frac{1}{n}\sum_{i=1}^n
\sum_w \Qbar_n(a,w)\ind(W_i=w).
\]
For each individual $i$, exactly one value of $w$ equals $W_i$, so the inner sum is
$\Qbar_n(a,W_i)$. Therefore,
\[
\sum_w \Qbar_n(a,w)P_n(W=w)
=
\frac{1}{n}\sum_{i=1}^n \Qbar_n(a,W_i).
\]

\subsection{By-Hand Example}

Suppose $W$ is binary and the sample summaries are as follows:
\[
\begin{array}{c|cc|c}
W & \Qbar_n(0,W) & \Qbar_n(1,W) & P_n(W)\\
\hline
0 & 2.680 & 2.875 & 0.370\\
1 & 2.222 & 2.327 & 0.630
\end{array}
\]
Then
\[
\psi_n(1)=2.875(0.370)+2.327(0.630)=2.530,
\]
and
\[
\psi_n(0)=2.680(0.370)+2.222(0.630)=2.392.
\]
Thus
\[
\widehat{\ATE}_{\mathrm{gcomp}}=2.530-2.392=0.138.
\]
The interpretation is that, after standardizing the treated and untreated subgroup
means to the same covariate distribution, the estimated mean outcome is about 0.138 units
higher under treatment than under no treatment.

\begin{estimationbox}
By-hand empirical g-computation:
\begin{enumerate}
\item List all observed covariate strata $w$.
\item Compute $\Qbar_n(1,w)$ and $\Qbar_n(0,w)$ within each stratum.
\item Compute $P_n(W=w)$ for each stratum.
\item Form $\psi_n(1)=\sum_w \Qbar_n(1,w)P_n(W=w)$.
\item Form $\psi_n(0)=\sum_w \Qbar_n(0,w)P_n(W=w)$.
\item Subtract to obtain $\widehat{\ATE}_{\mathrm{gcomp}}$.
\end{enumerate}
\end{estimationbox}
